% Л6 -- Ресурсы, файлы и асинхронность. Slides.
%
% Companion to the spoken script in ../L06-resources-and-async.md, which is
% the source of truth: the slides carry diagrams, code and short claims, and
% the delivery lives in the script only -- this deck carries no speaker
% notes. Slide text is in Russian because that is the content; everything
% else in this file is English.
%
% Build:  make          (latexmk -pdf; engine: pdflatex)
%
% Every listing was compiled and run on dotnet SDK 10.0.112 in a stock
% `dotnet new console`; outputs shown or described are real. See
% ../CLAUDE.md for the verification table.
%
% House rules: ../CLAUDE.md and ../../../CLAUDE.md. Theme: ../../../../theme/.

\documentclass[aspectratio=169,10pt,t]{beamer}

\usetheme{dotnet}
\usetikzlibrary{arrows.meta,positioning,calc,patterns}
\usepackage{array}   % ragged-right table columns

\dotnetlecture{Л6}

\title[Л6. Ресурсы, файлы, асинхронность]{Ресурсы, файлы и асинхронность}
\author{Дмитрий Курочкин}
\date{}

% Code sizes, one step down from \footnotesize. Smaller than 7pt does not
% carry to the back of the room.
\newcommand{\codesm}{\ttfamily\fontsize{7.6}{9.0}\selectfont}
\newcommand{\codexs}{\ttfamily\fontsize{7.0}{8.3}\selectfont}

\begin{document}

% ------------------------------------------------------------------ title --
\begin{frame}[plain]
  \titlepagednc{Лекция 6}{Ресурсы, файлы и асинхронность}{}{Дмитрий
    Курочкин}
\end{frame}

% ==================================================== 1. RESOURCE OWNERSHIP ==
\section{1. Владение ресурсом}

% 2
\begin{frame}[plain]
  \sectionopener{1}{Владение ресурсом}{За что отвечаете \hi{вы}, а не
    рантайм.}
\end{frame}

% 3
\begin{frame}{О чём первая часть}
  \slidetop
  \large
  \begin{cbul}{accent}
    \item managed и unmanaged ресурсы
    \item \texttt{IDisposable} и \texttt{using}
    \item три уровня работы с файлом
    \item пути к файлам
    \item финализаторы
  \end{cbul}
\end{frame}

% 4
\begin{frame}[plain]
  \vote{Файл открыли, прочитали и не закрыли. Объект стал недостижим, и
    сборщик мусора его убрал. Файл закрыт?}
\end{frame}

% 5
\begin{frame}{Два вида ресурсов}
  \slidetop
  \versus{managed}{%
    \begin{cbul}{accent}
      \item объекты в куче
      \item сборщик \hi{видит} их и убирает сам
    \end{cbul}}{unmanaged}{%
    \begin{cbul}{alt}
      \item файловый дескриптор, сокет, соединение с БД
      \item сборщик их \lo{не видит}
    \end{cbul}}

  \slidegap
  \aside{Для сборщика \texttt{FileStream} --- несколько десятков байт с числом
    внутри. Что число --- дескриптор, и ОС нужно сказать «я закончил», он не
    знает.}
\end{frame}

% 6
\begin{frame}[fragile]{\texttt{IDisposable}}
  \slidetop
  \claim{Договор: \hi{Dispose} значит «я закончил, освободи».}

\begin{lstlisting}[style=csharp,basicstyle=\ttfamily\small]
public interface IDisposable
{
    void Dispose();
}
\end{lstlisting}

  \slidegap
  {\renewcommand{\arraystretch}{1.35}\small
  \begin{tabular}{@{}p{3.6cm}p{7cm}@{}}
    \texttt{FileStream}  & закрывает дескриптор \\
    \texttt{StreamReader} & закрывает свой \texttt{FileStream} \\
    соединение с БД       & возвращается в пул \\
    сокет                 & закрывается \\
  \end{tabular}}
\end{frame}

% 7
\begin{frame}[fragile]{\texttt{using}: блок}
  \slidetop
  \claim{Старая форма: объект живёт до конца \hi{блока}.}

\begin{lstlisting}[style=csharp,basicstyle=\codesm]
string? FindHello(string path)
{
    using (StreamReader reader = new StreamReader(path))
    {
        string? line = reader.ReadLine();
        while (line != null && !line.Contains("Hello!"))
        {
            line = reader.ReadLine();
        }
        return line;
    }
}
\end{lstlisting}

  \slidegap
  \aside{Файл читается \hi{подряд}, с начала, до первого совпадения.}
\end{frame}

% 8
\begin{frame}[fragile]{\texttt{using}: объявление}
  \slidetop
  \claim{Современная форма: объект живёт до конца \hi{области}.}

\begin{lstlisting}[style=csharp,basicstyle=\codesm]
string? FindHello(string path)
{
    using StreamReader reader = new StreamReader(path);
    string? line = reader.ReadLine();
    while (line != null && !line.Contains("Hello!"))
    {
        line = reader.ReadLine();
    }
    return line;
}
\end{lstlisting}

  \slidegap
  \aside{Меньше отступов; несколько ресурсов подряд --- без лестницы блоков.}
\end{frame}

% 9
\begin{frame}[fragile]{Во что разворачивается \texttt{using}}
  \slidetop
  \claim{Обе формы --- это \hi{try / finally}: ресурс освобождён при любом
    исходе.}

\begin{lstlisting}[style=csharp,basicstyle=\codesm]
StreamReader reader = new StreamReader(path);
try
{
    // your code
}
finally
{
    if (reader != null)
    {
        reader.Dispose();
    }
}
\end{lstlisting}

  \slidegap
  \aside{Вернулись, бросили исключение, вышли через \texttt{break} ---
    \texttt{finally} выполнится.}
\end{frame}

% 10
\begin{frame}[fragile]{Несколько \texttt{using} подряд}
  \slidetop
  \claim{Освобождаются в \hi{обратном} порядке: последний созданный --- первым.}

\begin{lstlisting}[style=csharp,basicstyle=\codesm]
using FileStream stream = new FileStream(path, FileMode.Create);
using StreamWriter writer = new StreamWriter(stream);
writer.WriteLine("Hello!");
\end{lstlisting}

  \slidegap
  \begin{center}
    \begin{tikzpicture}[font=\small,
      st/.style={draw=hairline,line width=0.8pt,fill=codebg,
                 minimum width=5.6cm,minimum height=1.0cm,inner sep=4pt,
                 align=center},
      ar/.style={-{Straight Barb[length=1.8mm]},accent,line width=0.8pt}]
      \node[st,draw=accent] (w) at (0,0)
        {\textbf{1.} \texttt{writer.Dispose()}\\
         {\scriptsize сбрасывает буфер в поток}};
      \node[st] (s) at (6.8,0)
        {\textbf{2.} \texttt{stream.Dispose()}\\
         {\scriptsize закрывает дескриптор}};
      \draw[ar] (w) -- (s);
    \end{tikzpicture}
  \end{center}

  \slidegap
  \aside{Перепутанный порядок --- \lo{потерянный хвост} файла.}
\end{frame}

% 11
\begin{frame}[fragile]{Забытый \texttt{using}}
  \slidetop
  \claim{Незакрытый файл мешает вам \lo{самим}.}

\begin{lstlisting}[style=csharp,basicstyle=\codesm]
FileStream output = File.OpenWrite(path);
output.WriteByte(1);
string text = File.ReadAllText(path);   // throws
\end{lstlisting}

  \slidegap
\begin{lstlisting}[style=term,basicstyle=\ttfamily\fontsize{6.6}{8}\selectfont]
System.IO.IOException: The process cannot access the file
'/home/user/data.txt' because it is being used by another process.
\end{lstlisting}

  \slidegap
  \aside{«Другой процесс» --- это ваш дескриптор двумя строками выше. С \texttt{using}
    он освобождён \hi{детерминированно}, в \texttt{finally}.}
\end{frame}

% 12
\begin{frame}[fragile]{Свой тип: \texttt{LogWriter}}
  \vspace{0.3em}
\begin{lstlisting}[style=csharp,basicstyle=\ttfamily\fontsize{6.6}{7.6}\selectfont]
sealed class LogWriter(string path) : IDisposable
{
    private readonly StreamWriter _writer = new StreamWriter(path, append: true);
    private bool _disposed;

    public void Write(string line)
    {
        ObjectDisposedException.ThrowIf(_disposed, this);
        _writer.WriteLine($"{DateTime.UtcNow:O} {line}");
    }

    public void Dispose()
    {
        if (!_disposed)
        {
            _writer.Dispose();
            _disposed = true;
        }
    }
}
\end{lstlisting}

  \vspace{-0.2em}
  \aside{\hi{Повторный} Dispose безопасен; использование \hi{после} --- исключение.}
\end{frame}

% 13
\begin{frame}{Владеет --- освобождает}
  \slidetop
  \versus{создали сами}{%
    \begin{cbul}{accent}
      \item вы владелец
      \item освобождаете: \texttt{using} или \texttt{Dispose}
    \end{cbul}}{получили снаружи}{%
    \begin{cbul}{alt}
      \item вы не владелец
      \item не освобождаете: это сделает создатель
    \end{cbul}}

  \slidegap
  \hrulethin{0.42}
  \rulegap
  {\large Владелец \hi{один}: двое --- двойное освобождение, никто --- утечка.}
\end{frame}

% 14
\begin{frame}[fragile]{Уровень 1: \texttt{File}}
  \slidetop
  \claim{Одна строка: открыл, прочитал, закрыл --- \hi{сам}.}

\begin{lstlisting}[style=csharp,basicstyle=\ttfamily\small]
string text = File.ReadAllText(path);
string[] lines = File.ReadAllLines(path);
File.WriteAllText(path, text);
\end{lstlisting}

  \slidegap
  \aside{Файл целиком в памяти. Верхнюю границу размера назвать не можете ---
    \lo{не читаете целиком}.}
\end{frame}

% 15
\begin{frame}{Уровень 2: \texttt{StreamReader}}
  \slidetop
  \claim{Читает блоками, отдаёт строки: в памяти один \hi{блок} и одна
    \hi{строка}.}

  \begin{center}
    \begin{tikzpicture}[font=\small,
      cell/.style={draw=hairline,line width=0.8pt,fill=codebg,
                   minimum width=1.2cm,minimum height=0.8cm,inner sep=0pt},
      ar/.style={-{Straight Barb[length=1.8mm]},accent,line width=0.7pt}]
      \foreach \i in {0,...,7} \node[cell] (b\i) at (\i*1.2,0) {};
      \node[cell,draw=accent,line width=1.2pt] at (b3) {};
      \node[muted,font=\scriptsize,above=0.3em] at (b3.north) {файл на диске};
      \node[draw=accent,line width=0.9pt,fill=codebg,minimum width=2.6cm,
        minimum height=0.8cm] (buf) at (3.6,-1.9) {буфер: блок};
      \node[draw=accent,line width=0.9pt,fill=codebg,minimum width=2.2cm,
        minimum height=0.8cm] (ln) at (7.4,-1.9) {строка};
      \draw[ar] (b3.south) -- (buf.north);
      \draw[ar] (buf.east) -- node[above,font=\scriptsize,text=accent]
        {\texttt{ReadLine}} (ln.west);
    \end{tikzpicture}
  \end{center}

  \slidegap
  \aside{Гигабайтный файл --- в постоянной памяти. Первый уровень, где нужен
    \texttt{using}.}
\end{frame}

% 16
\begin{frame}[fragile]{Уровень 3: \texttt{FileStream}}
  \slidetop
  \claim{Байты, позиция --- и \hi{доступ к файлу} задаётся явно.}

  \begin{minipage}[t]{0.50\textwidth}
\begin{lstlisting}[style=csharp,basicstyle=\codesm]
FileStream stream = new FileStream(
    path,
    FileMode.Open,
    FileAccess.Read,
    FileShare.Read);
\end{lstlisting}
  \end{minipage}\hfill
  \begin{minipage}[t]{0.46\textwidth}
    {\renewcommand{\arraystretch}{1.5}\small
    \begin{tabular}{@{}p{2.2cm}p{3.8cm}@{}}
      \hi{FileMode}   & создать, открыть, дописать \\
      \hi{FileAccess} & чтение, запись \\
      \hi{FileShare}  & что разрешено другим \\
    \end{tabular}}
  \end{minipage}

  \slidegap
  \aside{Пока файл открыт у вас, другие смогут его читать, но не менять. Это
    контракт с ОС --- выбирайте его сознательно.}
\end{frame}

% 17
\begin{frame}[fragile]{Путь: разделители}
  \slidetop
  \claim{Путь собирает \hi{Path}, а не конкатенация строк.}

\begin{lstlisting}[style=csharp,basicstyle=\ttfamily\small]
Path.Combine("data", "logs", "app.txt")
\end{lstlisting}

  \slidegap
  \versus{Windows}{\texttt{data\textbackslash logs\textbackslash app.txt}}{%
    Linux, macOS}{\texttt{data/logs/app.txt}}

  \slidegap
  \aside{Разделитель подставляется по платформе. Склеенный строкой путь
    работает на \lo{одной}.}
\end{frame}

% 18
\begin{frame}[fragile]{\texttt{Path}: собрать и разобрать}
  \slidetop
  \claim{Не режьте путь руками: эти методы знают про все платформы.}

\begin{lstlisting}[style=csharp,basicstyle=\codesm]
string Path.Combine(string path1, string path2)
string Path.Join(string path1, string path2)
string Path.GetExtension(string path)
string Path.GetFileName(string path)
string Path.GetDirectoryName(string path)
\end{lstlisting}

  \slidegap
  \aside{«Если файл есть, открыть» --- гонка между проверкой и открытием.
    Открывайте и обрабатывайте \hi{исключение}.}
\end{frame}

% 19
\begin{frame}{Относительный и абсолютный путь}
  \slidetop
  \versus{абсолютный}{%
    \begin{cbul}{accent}
      \item не зависит от места запуска
      \item рядом с приложением --- от папки приложения
      \item внешний файл --- путь из конфигурации
    \end{cbul}}{относительный}{%
    \begin{cbul}{alt}
      \item отсчёт от \lo{рабочей директории процесса}, не от сборки
      \item \texttt{dotnet run}, публикация, контейнер --- три разных места
    \end{cbul}}

  \slidegap
  \hrulethin{0.42}
  \rulegap
  {\large Регистр: Windows не различает, Linux \lo{различает} ---
    \texttt{Config.json} и \texttt{config.json} там два файла.}
\end{frame}

% 20
\begin{frame}{Где приложению можно работать с файлами}
  \slidetop
  \claim{Приложение работает только со \hi{своей} папкой.}

  \versus{своя папка --- читаем}{%
    \begin{cbul}{accent}
      \item конфигурация, шаблоны, данные приложения
      \item в контейнере может быть только для чтения
    \end{cbul}}{временная --- пишем}{%
    \begin{cbul}{alt}
      \item всё, что нужно записать
      \item \texttt{Path.GetTempPath()}
    \end{cbul}}
\end{frame}

% 21
\begin{frame}[fragile]{Финализатор}
  \slidetop
  \claim{Метод с тильдой: сборщик вызовет его перед тем, как убрать объект.}

\begin{lstlisting}[style=csharp,basicstyle=\ttfamily\small]
class Holder
{
    ~Holder()
    {
        // called by the GC before the object is removed
    }
}
\end{lstlisting}

  \slidegap
  \aside{Так \texttt{FileStream} закрывается «когда-нибудь». Когда именно ---
    решает сборщик.}
\end{frame}

% 22
\begin{frame}{Три причины не писать финализатор}
  \slidetop
  {\renewcommand{\arraystretch}{1.9}
  \begin{tabular}{@{}p{4.2cm}p{8.6cm}@{}}
    \lo{\textbf{недетерминированность}} & через миллисекунду, через час или
      никогда \\
    \lo{\textbf{цена}} & объект переживает лишнюю сборку --- дороже обычного
      всегда \\
    \lo{\textbf{сложность}} & отдельный поток, неизвестный момент, другие
      объекты уже могли быть собраны \\
  \end{tabular}}

  \slidegap
  \aside{Вместо финализатора --- \hi{Dispose}. Сырой дескриптор ОС --- готовая
    обёртка из стандартной библиотеки.}
\end{frame}

% 23
\begin{frame}[plain]
  \keyline{Сборщик убирает память, а не ресурсы.\\ Ресурс освобождает
    владелец: один раз, через \texttt{using}.}
\end{frame}

% ================================================================ 2. ASYNC ==
\section{2. Асинхронность}

% 24
\begin{frame}[plain]
  \sectionopener{2}{Асинхронность}{Поток --- \hi{дорогой} ресурс, и он не
    должен ждать.}
\end{frame}

% 25
\begin{frame}[plain]
  \vote{Запрос идёт в базу данных, ответ приходит через 50 мс. Что делает
    поток, обрабатывающий запрос, эти 50 мс?}
\end{frame}

% 26
\begin{frame}{Что делает поток, пока ждёт}
  \slidetop
  \claim{Ждёт: занимает стек и запись в ядре, и не делает \lo{ничего}.}

  \begin{center}
    \begin{tikzpicture}[font=\small]
      \node[draw=accent,line width=1.0pt,fill=codebg,minimum width=1.6cm,
        minimum height=1.0cm] (a) at (0,0) {};
      \node[draw=alt,line width=1.0pt,dashed,fill=paper,minimum width=7.4cm,
        minimum height=1.0cm,anchor=west] (w) at (a.east) {};
      \node[draw=accent,line width=1.0pt,fill=codebg,minimum width=1.6cm,
        minimum height=1.0cm,anchor=west] (b) at (w.east) {};
      \node[font=\scriptsize,text=accent,below=0.3em] at (a.south) {принял};
      \node[font=\scriptsize,text=alt,below=0.3em] at (w.south)
        {база думает: 50 мс, поток заблокирован};
      \node[font=\scriptsize,text=accent,below=0.3em] at (b.south) {собрал};
      \draw[-{Straight Barb[length=1.8mm]},muted,line width=0.6pt]
        ([yshift=-1.3cm]a.west) -- node[below,font=\scriptsize,text=muted]
        {время} ([yshift=-1.3cm]b.east);
    \end{tikzpicture}
  \end{center}

  \slidegap
  \aside{За эти 50 мс процессор мог выполнить десятки миллионов инструкций.
    Поток не выполнил ни одной.}
\end{frame}

% 27
\begin{frame}{Поток на запрос и асинхронность}
  \slidetop
  \begin{center}
    \begin{tikzpicture}[font=\small]
      \foreach \r in {0,...,4}
        \foreach \c in {0,...,9}
          \node[draw=alt,line width=0.6pt,fill=codebg,minimum size=0.28cm,
            inner sep=0pt] at (\c*0.36,-\r*0.36) {};
      \node[alt,font=\footnotesize\bfseries] at (1.6,-2.2) {1000 потоков, все ждут};

      \foreach \c in {0,...,7}
        \node[draw=accent,line width=0.9pt,fill=accent!15!paper,minimum size=0.42cm,
          inner sep=0pt] at (7.3+\c*0.5,0) {};
      \foreach \x/\y in {7.4/-0.9,8.0/-1.3,8.7/-0.8,9.3/-1.4,9.9/-0.9,10.5/-1.3,
                         7.7/-1.8,8.4/-1.7,9.0/-2.0,9.7/-1.8,10.3/-2.0}
        \node[fill=hairline,circle,minimum size=0.18cm,inner sep=0pt] at (\x,\y) {};
      \node[accent,font=\footnotesize\bfseries] at (9.2,-2.6)
        {8 потоков работают, ожидания --- объекты};
    \end{tikzpicture}
  \end{center}

  \slidegap
  \aside{Один запрос не станет быстрее: те же 50 мс. Меняется то, сколько
    запросов сервер держит одновременно.}
\end{frame}

% 28
\begin{frame}{Асинхронность и параллелизм}
  \slidetop
  {\renewcommand{\arraystretch}{1.9}
  \begin{tabular}{@{}p{2.6cm}p{6.6cm}p{2.4cm}@{}}
    \colhead{muted}{}      & \colhead{muted}{потоки}   & \colhead{muted}{когда} \\
    \textbf{синхронно}     & один, ждёт                 & \\
    \hi{\textbf{асинхронно}} & мало, никто не ждёт      & \hi{ввод-вывод} \\
    \lo{\textbf{параллельно}} & по числу ядер, все считают & \lo{вычисления} \\
  \end{tabular}}

  \slidegap
  \aside{Асинхронность про \hi{ожидание}, параллелизм про \lo{вычисления}.
    Лишних ядер асинхронность не даёт.}
\end{frame}

% 29
\begin{frame}{Правило}
  \slidetop
  \versus{ввод-вывод}{%
    {\large \hi{асинхронно}, через \texttt{await}}}{вычисления}{%
    {\large\dncnohyph \lo{параллельно}, если можно разрезать; синхронно, если нельзя}}

  \slidegap
  \hrulethin{0.42}
  \rulegap
  \aside{Ввод-вывод в \texttt{Task.Run} --- поток из пула, который только ждёт.
    Тяжёлое вычисление под \texttt{async} --- всё тот же занятый поток.}
\end{frame}

% 30
\begin{frame}[plain]
  \vote{Что такое \texttt{Task}?}
\end{frame}

% 31
\begin{frame}{\texttt{Task} --- обещание}
  \slidetop
  \claim{Не поток и не операция: результат \hi{будет}, возможно, не сейчас.}

  \begin{center}
    \begin{tikzpicture}[font=\small,
      row/.style={draw=hairline,line width=0.6pt,fill=codebg,
                  minimum width=5.4cm,minimum height=0.8cm,inner sep=4pt},
      src/.style={draw=hairline,line width=0.8pt,fill=paper,
                  minimum width=4.4cm,minimum height=0.9cm,inner sep=3pt,
                  align=center},
      ar/.style={-{Straight Barb[length=1.8mm]},line width=0.8pt}]
      \node[row,draw=accent,line width=1.0pt] (s) at (7.4,0.85) {состояние};
      \node[row] (r) at (7.4,0)    {результат / исключение};
      \node[row] (c) at (7.4,-0.85) {список continuations};
      \node[src] (run) at (1.4,0.7)  {\texttt{Task.Run}\\
        {\scriptsize поток из пула}};
      \node[src] (io)  at (1.4,-0.7) {ввод-вывод\\
        {\scriptsize ОС сообщит}};
      \draw[ar,accent] (run.east) -- ([xshift=-2pt]s.west |- run.east);
      \draw[ar,alt] (io.east) -- ([xshift=-2pt]c.west |- io.east);
    \end{tikzpicture}
  \end{center}
\end{frame}

% 32
\begin{frame}[fragile]{Источник 1: \texttt{Task.Run}}
  \slidetop
  \claim{Вычисление на потоке из пула: за задачей \hi{стоит поток}.}

\begin{lstlisting}[style=csharp,basicstyle=\ttfamily\small]
Task<int> task = Task.Run(() => Compute());
\end{lstlisting}

  \slidegap
  \aside{Инструмент для вычислений. Для ввода-вывода не нужен: поток из пула
    будет только ждать.}
\end{frame}

% 33
\begin{frame}[fragile]{Источник 2: ввод-вывод}
  \slidetop
  \claim{Асинхронное чтение: за задачей \hi{потока нет}.}

\begin{lstlisting}[style=csharp,basicstyle=\ttfamily\small]
Task<string> task = File.ReadAllTextAsync(path);
\end{lstlisting}

  \slidegap
  \aside{Библиотека попросила ОС прочитать и сообщить. Между запросом и ответом
    задача есть, потока --- нет.}
\end{frame}

% 34
\begin{frame}[fragile]{Пауза: два способа подождать секунду}
  \slidetop
  \versus{\texttt{Task.Delay}}{%
    \begin{cbul}{accent}
      \item регистрирует таймер в ОС
      \item поток \hi{свободен}
      \item таймер завершит задачу и запустит продолжение
    \end{cbul}}{\texttt{Thread.Sleep}}{%
    \begin{cbul}{alt}
      \item блокирует поток
      \item поток ничего не делает и никому не достаётся
      \item в сервисе --- \lo{никогда}
    \end{cbul}}
\end{frame}

% 35
\begin{frame}[plain]
  \statement{У асинхронного метода нет своего потока.}
\end{frame}

% 36
\begin{frame}[fragile]{Пять \texttt{await} подряд}
  \slidetop
  \claim{Каждый \hi{await} --- место, где метод может улететь в ожидание.}

\begin{lstlisting}[style=csharp,basicstyle=\codesm]
async Task<string> HandleAsync(int id)
{
    User user = await LoadUserAsync(id);              // 1
    int[] orders = await LoadOrdersAsync(user);       // 2
    decimal price = await CalcPriceAsync(orders);     // 3
    await SaveAsync(price);                           // 4
    await NotifyAsync(user);                          // 5
    return user.Name;
}
\end{lstlisting}

  \slidegap
  \aside{Выглядит как пять строк подряд, а по смыслу это callback: «всё, что
    ниже, выполни, когда это завершится».}
\end{frame}

% 37
\begin{frame}{Цепочка состояний}
  \slidetop
  \claim{Пять \hi{await} режут метод на шесть кусков; номер куска ---
    состояние.}

  \begin{center}
    \begin{tikzpicture}[font=\small,
      pc/.style={draw=hairline,line width=0.8pt,fill=codebg,
                 minimum width=1.4cm,minimum height=1.0cm,inner sep=0pt},
      ar/.style={-{Straight Barb[length=1.6mm]},alt,line width=0.8pt}]
      \foreach \i in {0,...,5} \node[pc] (p\i) at (\i*2.2,0) {\textbf{\i}};
      \foreach \i/\j in {0/1,1/2,2/3,3/4,4/5} {
        \draw[ar] (p\i.east) -- node[above,font=\scriptsize\ttfamily,text=alt]
          {await} (p\j.west);
      }
      \node[muted,font=\scriptsize,below=0.5em] at (p2.south)
        {состояние = номер куска};
    \end{tikzpicture}
  \end{center}

  \slidegap
  \aside{На каждом \texttt{await} состояние запоминается, метод улетает; когда
    ожидание кончилось, продолжение читает номер и прыгает на нужный кусок.}
\end{frame}

% 38
\begin{frame}[fragile]{Исключение в асинхронном методе}
  \slidetop
  \claim{Исключение лежит в задаче и вылетает в строке \hi{await}.}

\begin{lstlisting}[style=csharp,basicstyle=\ttfamily\fontsize{6.8}{7.9}\selectfont]
async Task FailAsync()
{
    await Task.Delay(10);
    throw new InvalidOperationException("boom");
}

try
{
    await FailAsync();   // the exception is thrown here
}
catch (InvalidOperationException ex)
{
    Console.WriteLine(ex.Message);
}
\end{lstlisting}

  \vspace{0.2em}
  \aside{Стека вызова уже нет, поэтому исключение ждёт в задаче.}
\end{frame}

% 39
\begin{frame}{Как работает переключение: четыре шага}
  \slidetop
  {\renewcommand{\arraystretch}{1.55}\small
  \begin{tabular}{@{}p{0.7cm}p{2.8cm}p{8.8cm}@{}}
    \hi{\textbf{1}} & компилятор & метод становится \hi{state machine}: локальные
      переменные --- поля, номер состояния, \texttt{MoveNext} \\
    \hi{\textbf{2}} & \texttt{await} & завершена? да --- идём дальше; нет ---
      регистрируем продолжение и \hi{возвращаемся} \\
    \hi{\textbf{3}} & ОС & операция завершилась --- продолжение ставится в
      очередь thread pool \\
    \hi{\textbf{4}} & поток пула & \texttt{MoveNext} продолжает с
      сохранённого состояния \\
  \end{tabular}}

  \slidegap
  \aside{Часто задача уже завершена: тогда никакого переключения, метод идёт
    дальше на том же потоке.}
\end{frame}

% 40
\begin{frame}{Шаг 3: куда девается поток}
  \slidetop
  \begin{center}
    \begin{tikzpicture}[font=\footnotesize,
      bx/.style={draw=hairline,line width=0.8pt,fill=codebg,
                 minimum width=2.6cm,minimum height=0.95cm,inner sep=3pt,
                 align=center},
      ar/.style={-{Straight Barb[length=1.6mm]},line width=0.7pt}]
      \node[bx] (t) at (0,1.2) {ваш поток:\\ «прочитай сокет»};
      \node[bx] (os1) at (3.6,1.2) {ОС: запомнила\\ запрос};
      \node[bx,draw=accent] (pool) at (7.2,1.2) {поток --- в пул,\\ берёт другую работу};
      \draw[ar,muted] (t) -- (os1);
      \draw[ar,accent] (os1.east) -- (pool.west);

      \node[alt,font=\scriptsize] at (5.4,0.1)
        {ожидание: записи в таблице ядра, потока нет};

      \node[bx] (os2) at (0,-1.2) {ОС: данные\\ пришли};
      \node[bx] (svc) at (3.0,-1.2) {служебный поток:\\ задаче --- результат};
      \node[bx] (q) at (6.0,-1.2) {очередь\\ thread pool};
      \node[bx,draw=accent] (w) at (9.0,-1.2) {любой поток пула:\\ \texttt{MoveNext}};
      \draw[ar,muted] (os2) -- (svc);
      \draw[ar,muted] (svc) -- (q);
      \draw[ar,accent] (q) -- (w);
    \end{tikzpicture}
  \end{center}

  \slidegap
  \aside{«There is no thread»: между запросом и ответом вашим запросом не занят
    ни один поток.}
\end{frame}

% 41
\begin{frame}[plain]
  \vote{Код до \texttt{await} выполнялся на потоке номер 5 из пула. На каком
    потоке выполнится код после \texttt{await}?}
\end{frame}

% 42
\begin{frame}{Куски метода на разных потоках}
  \slidetop
  \claim{Метод начался на одном потоке и закончится на \hi{другом}.}

  \begin{center}
    \begin{tikzpicture}[font=\small,
      lane/.style={draw=none},
      pc/.style={draw=accent,line width=1.0pt,fill=codebg,minimum width=2.2cm,
                 minimum height=0.8cm,inner sep=0pt},
      ar/.style={-{Straight Barb[length=1.6mm]},alt,line width=0.7pt,dashed}]
      \foreach \y/\n in {1.5/5, 0/2, -1.5/7} {
        \node[muted,font=\scriptsize,anchor=east] at (-0.3,\y) {поток \n};
        \draw[hairline,line width=0.5pt] (0,\y) -- (11.4,\y);
      }
      \node[pc] (a) at (1.4,1.5)  {кусок 0};
      \node[pc] (b) at (5.6,0)    {кусок 1};
      \node[pc] (c) at (9.8,-1.5) {кусок 2};
      \draw[ar] (a.east) -- node[above,font=\scriptsize,text=alt] {await} (b.north west);
      \draw[ar] (b.east) -- node[above,font=\scriptsize,text=alt] {await} (c.north west);
    \end{tikzpicture}
  \end{center}

  \slidegap
  \aside{Пятый поток давно обслуживает другие запросы. Продолжение берёт тот,
    кто свободен.}
\end{frame}

% 43
\begin{frame}{Что переживает \texttt{await}}
  \slidetop
  \versus{переживает}{%
    \begin{cbul}{accent}
      \item локальные переменные: они поля state machine в куче
      \item \texttt{this}
      \item \texttt{AsyncLocal}
    \end{cbul}}{не переживает}{%
    \begin{cbul}{alt}
      \item данные, локальные для потока
      \item блокировка \texttt{lock}: \texttt{await} внутри запрещён
    \end{cbul}}

  \slidegap
  \aside{Большой массив в локальной переменной живёт всё время ожидания, даже
    если после \texttt{await} он не нужен.}
\end{frame}

% 44
\begin{frame}{Thread pool}
  \slidetop
  \claim{Очередь элементов работы и потоки, которые её разбирают.}

  \begin{center}
    \begin{tikzpicture}[font=\footnotesize,
      it/.style={draw=hairline,line width=0.8pt,fill=codebg,
                 minimum width=2.1cm,minimum height=0.7cm,inner sep=2pt},
      th/.style={draw=accent,line width=1.0pt,fill=accent!12!paper,
                 minimum width=1.3cm,minimum height=0.7cm,inner sep=2pt},
      ar/.style={-{Straight Barb[length=1.6mm]},muted,line width=0.7pt}]
      \node[it] (q1) at (0,0.8)  {продолжение};
      \node[it] (q2) at (0,0)    {\texttt{Task.Run}};
      \node[it] (q3) at (0,-0.8) {таймер};
      \node[muted,font=\scriptsize] at (0,1.4) {очередь};
      \foreach \i in {0,1,2,3} \node[th] (w\i) at (5.4,1.2-\i*0.8) {поток};
      \node[muted,font=\scriptsize] at (5.4,1.9) {по числу ядер};
      \draw[ar] (q2.east) -- ([xshift=-2pt]w1.west);
      \draw[ar] (q1.east) -- ([xshift=-2pt]w0.west);
      \draw[ar] (q3.east) -- ([xshift=-2pt]w2.west);
    \end{tikzpicture}
  \end{center}

  \vspace{0.4em}
  \aside{Если очередь растёт, пул добавляет потоки, но \lo{медленно}: он не
    отличает тяжёлую работу от заблокированных потоков.}
\end{frame}

% 45
\begin{frame}{Голодание пула}
  \slidetop
  \begin{center}
    \begin{tikzpicture}[font=\footnotesize,
      it/.style={draw=hairline,line width=0.8pt,fill=codebg,
                 minimum width=0.9cm,minimum height=0.6cm,inner sep=2pt},
      th/.style={draw=alt,line width=1.0pt,fill=alt!12!paper,
                 minimum width=1.4cm,minimum height=0.6cm,inner sep=2pt}]
      \foreach \i in {0,...,7} \node[it] at (\i*1.0,0) {};
      \node[muted,font=\scriptsize,above=0.5em] at (3.5,0.3)
        {очередь растёт: продолжениям некому выполниться};
      \foreach \i in {0,...,3} \node[th] at (9.6+\i*1.6-3.2,-1.5) {ждёт};
      \node[alt,font=\scriptsize] at (9.6,-2.15)
        {все потоки заблокированы синхронным ожиданием};
    \end{tikzpicture}
  \end{center}

  \slidegap
  \begin{center}
    {\LARGE\bfseries\lo{Не блокируй поток!}}
  \end{center}

  \aside{Симптом: задержки в секунды при простаивающем процессоре и быстрой
    базе; очередь пула растёт.}
\end{frame}

% 46
\begin{frame}{Один запрос от начала до конца}
  \slidetop
  \begin{center}
    \begin{tikzpicture}[font=\footnotesize,
      bx/.style={draw=hairline,line width=0.8pt,fill=codebg,
                 minimum height=1.05cm,inner sep=4pt,align=center},
      ar/.style={-{Straight Barb[length=1.6mm]},line width=0.7pt}]
      \foreach \y/\n in {1.5/{поток 3}, 0/{база}, -1.5/{поток 7}} {
        \node[muted,font=\scriptsize,anchor=east] at (0.9,\y) {\n};
        \draw[hairline,line width=0.5pt] (1.0,\y) -- (13.0,\y);
      }
      \node[bx,minimum width=4.5cm,draw=accent] (a1) at (3.4,1.5)
        {\textbf{1. Запрос А идёт вниз}\\
         \textcolor{muted}{\scriptsize до драйвера БД; на каждом \texttt{await}}\\
         \textcolor{muted}{\scriptsize возвращается вверх}};
      \node[bx,minimum width=3.3cm] (b) at (7.7,1.5)
        {\textbf{2. Поток 3 свободен}\\
         \textcolor{muted}{\scriptsize берёт запрос Б}};
      \node[bx,minimum width=4.2cm,draw=alt,dashed,fill=paper] (db) at (8.3,0)
        {\textbf{3. База отвечает через 40 мс}\\
         \textcolor{muted}{\scriptsize ни один поток не занят}};
      \node[bx,minimum width=4.5cm,draw=accent] (a2) at (10.5,-1.5)
        {\textbf{4. Запрос А идёт вверх}\\
         \textcolor{muted}{\scriptsize на потоке 7, до ответа клиенту}};
      \draw[ar,muted] (a1.south) |- (db.west);
      \draw[ar,accent] (db.east) -| (a2.north);
    \end{tikzpicture}
  \end{center}

  \slidegap
  \aside{Запрос А начался на потоке 3, закончился на потоке 7 и в
    промежутке не занимал \hi{ни одного}.}
\end{frame}

% 47
\begin{frame}[plain]
  \keyline{Await --- не ожидание, а регистрация продолжения и возврат.\\
    Между кусками метода потока нет вообще.}
\end{frame}

% 48
\begin{frame}[plain]
  \vote{Два независимых похода по 50 мс, два \texttt{await} подряд. Сколько
    займёт метод?}
\end{frame}

% 49
\begin{frame}[fragile]{\texttt{Task.WhenAll}}
  \slidetop
  \claim{Задачи созданы \hi{до} await; await один --- на \texttt{WhenAll}.}

\begin{lstlisting}[style=csharp,basicstyle=\codesm]
async Task<int> TotalLengthAsync(string first, string second, CancellationToken ct)
{
    Task<string> firstTask = File.ReadAllTextAsync(first, ct);
    Task<string> secondTask = File.ReadAllTextAsync(second, ct);

    string[] texts = await Task.WhenAll(firstTask, secondTask);

    return texts[0].Length + texts[1].Length;
}
\end{lstlisting}

  \slidegap
  \aside{50 мс вместо 100, ни одного лишнего потока. Но \lo{десять тысяч}
    задач --- это десять тысяч одновременных запросов: параллелизм ограничивают.}
\end{frame}

% 50
\begin{frame}[fragile]{\texttt{CancellationToken}}
  \vspace{0.3em}
  \claim{Последний параметр каждого метода --- и передаётся \hi{дальше}.}

\begin{lstlisting}[style=csharp,basicstyle=\ttfamily\fontsize{6.6}{7.6}\selectfont]
async Task<string> HandleAsync(string path, CancellationToken ct)
{
    int count = await CountOrdersAsync(path, ct);
    return "orders: " + count;
}

async Task<int> CountOrdersAsync(string path, CancellationToken ct)
{
    string[] lines = await ReadLinesAsync(path, ct);
    return lines.Length;
}

Task<string[]> ReadLinesAsync(string path, CancellationToken ct)
{
    return File.ReadAllLinesAsync(path, ct);
}
\end{lstlisting}

  \vspace{-0.3em}
  \aside{Клиент закрыл вкладку --- отмена доходит до самого нижнего вызова.}
\end{frame}

% 51
\begin{frame}[fragile]{Токен в своём цикле}
  \slidetop
  \claim{Ввод-вывод прервётся сам; свой цикл токен нужно \hi{спросить}.}

\begin{lstlisting}[style=csharp,basicstyle=\ttfamily\fontsize{6.8}{7.9}\selectfont]
int Sum(int[] items, CancellationToken ct)
{
    int total = 0;
    for (int i = 0; i < items.Length; i++)
    {
        ct.ThrowIfCancellationRequested();   // throws OperationCanceledException
        total += items[i];
    }
    if (ct.IsCancellationRequested)          // asks without throwing
    {
        return -1;
    }
    return total;
}
\end{lstlisting}

  \vspace{0.2em}
  \aside{Отмена --- исключение, его \lo{не глушат}. Код, который токен не
    проверяет, доработает до конца.}
\end{frame}

% 52
\begin{frame}[fragile]{Файлы, асинхронно}
  \slidetop
  \begin{minipage}[t]{0.67\textwidth}
\begin{lstlisting}[style=csharp,basicstyle=\codexs]
async Task CopyNonEmptyLinesAsync(
    string source, string target, CancellationToken ct)
{
    using StreamReader reader = new StreamReader(source);
    await using StreamWriter writer = new StreamWriter(target);

    string? line = await reader.ReadLineAsync(ct);
    while (line != null)
    {
        if (line.Length > 0)
        {
            await writer.WriteLineAsync(line.AsMemory(), ct);
        }
        line = await reader.ReadLineAsync(ct);
    }
}
\end{lstlisting}
  \end{minipage}\hfill
  \begin{minipage}[t]{0.29\textwidth}
    \eyebrow{что здесь используется}\par\vspace{0.6em}
    \small
    \begin{cbul}{accent}
      \item \texttt{using} --- читателю нечего сбрасывать
      \item \texttt{await using} --- писатель сбрасывает буфер
      \item \texttt{ReadLineAsync(ct)}
      \item \texttt{AsMemory()}: \texttt{Memory}, не \texttt{Span}
    \end{cbul}
  \end{minipage}
\end{frame}

% 53
\begin{frame}[fragile]{\texttt{IAsyncEnumerable}}
  \slidetop
\begin{lstlisting}[style=csharp,basicstyle=\codexs]
async IAsyncEnumerable<string> ErrorsAsync(
    string path, [EnumeratorCancellation] CancellationToken ct = default)
{
    await foreach (string line in File.ReadLinesAsync(path, ct))
    {
        if (line.Contains("ERROR"))
        {
            yield return line;
        }
    }
}

await foreach (string error in ErrorsAsync("app.log", token))
{
    Console.WriteLine(error);
}
\end{lstlisting}

  \vspace{0.3em}
  \aside{Строки не загружены заранее: в памяти одна. Атрибут на токене нужен,
    чтобы токен из \texttt{await foreach} доехал до тела метода.}
\end{frame}

% 54
\begin{frame}{Три опасности}
  \slidetop
  {\renewcommand{\arraystretch}{1.9}
  \begin{tabular}{@{}p{4.0cm}p{8.8cm}@{}}
    \lo{\textbf{\texttt{.Result} / \texttt{.Wait()}}} & поток выбыл ---
      голодание пула \\
    \lo{\textbf{\texttt{async void}}} & исключению некуда деться: процесс
      падает \\
    \lo{\textbf{забытый \texttt{await}}} & исключение исчезает, порядок
      нарушен \\
  \end{tabular}}

  \slidegap
  \aside{Асинхронность \hi{до конца}: если внизу есть \texttt{await}, он есть
    во всех вызывающих. \texttt{async void} --- только для обработчиков
    событий.}
\end{frame}

% 55
\begin{frame}[plain]
  \keyline{Await регистрирует продолжение и отдаёт поток.\\ Всё, что заставляет
    поток ждать, ломает эту схему.}
\end{frame}

\end{document}
